\section{Bounded growth ideals}
\label{sec:analytic-bounded-growth-ideals}
Integrability permits a finite loss of vanishing order. Bounded growth records the order without that loss. For analytic singularities its coherence can be proved using a finite normalization, rather than a general coherent direct-image theorem.

\begin{definition}\label{def:analytic-bounded-growth-ideal}
For a quasi-plurisubharmonic function $v$, let $\mathcal I_\infty(v)$ be the sheaf of holomorphic germs $f$ satisfying
\[
 \log|f|^2\leq v+O(1),
 \quad\text{equivalently}\quad |f|^2\leq C\exp(v)
\]
on a neighborhood of their base point, for some finite $C$. The zero germ belongs to this sheaf. Bounded changes of $v$ leave it unchanged. The triangle inequality and local boundedness of holomorphic coefficients make it a sheaf of ideals. It is contained in $\mathcal I(v)$, since its weighted integrand is locally bounded.
\end{definition}
\index{ideal!bounded growth}\index{$\mathcal I_\infty$}

\begin{theorem}\label{thm:analytic-bounded-growth-coherence}
If $v$ has analytic singularities, the sheaf $\mathcal I_\infty(v)$ is coherent, without a compactness assumption or a specified global ideal type. More precisely, for a nonzero holomorphic tuple $g=(g_1,\ldots,g_r)$ on a coordinate neighborhood $U$ and positive integers $p,q$, the sheaf
\[
 \mathcal J_{p/q}(g)=\{f\in\mathcal O_U:
               |f|\leq C|g|^{p/q}\text{ locally for some }C\}
\]
is a coherent ideal. Here $|g|=(\sum_j|g_j|^2)^{1/2}$.
\end{theorem}
\begin{proof}
Coherence is local. Shrink $U$ around a chosen point $x$ so that it is connected, and choose $g_i$ not identically zero. Put $U_i=\{g_i\ne0\}$. This open set is connected: the bounded holomorphic removal theorem extends every locally constant function with values zero and one to $U$, where the identity theorem makes it constant. We first construct a finite analytic space that records the fractional power.

Blow up the ideal $(g_1,\ldots,g_r)$ using \cref{prop:analytic-blowup-universal}, and denote the result by $B\to U$. A point of $B\subset U\times\mathbb P^{r-1}$ specifies a line in $\mathbb C^r$. The total space of those lines maps to $U\times\mathbb C^r$, with coordinates $(z,a)$, by retaining the vector on that line. This map is proper: the inverse image of a compact set is a closed subset of the product of its compact inverse image in $B$ with a bounded set of vectors. Its reduced image $C$ is analytic by \cref{thm:proper-analytic-image}. It is invariant under scalar multiplication of $a$. On $\{g_i\ne0\}$ it is the smooth line family
\[
 a_j=\tau g_j(z),\qquad \tau\in\mathbb C.
\]
This family is dense in $C$, since the corresponding open subset is dense in $B$ and hence in its line bundle. In particular $C$ is irreducible after shrinking the connected base; the dense open family is irreducible by the identity theorem.

In $C\times\mathbb C_b$ impose
\[
 b^q=a_i g_i^{q-1}.
\]
Retain the closure $C'$ of its part over $\{g_i\ne0\}$. It is analytic by finite local branch decomposition, and finite over $C$, since the displayed monic equation bounds $b$ over compact sets and has finite fibers. Its dense open subset is parametrized by
\[
 (z,t)\longmapsto
       (z,a_j=t^qg_j(z),b=tg_i(z)).
\]
It is irreducible. Normalize $C'$ by \cref{thm:analytic-normalization}, obtaining a finite map $\nu:Y\to C$. Over $U_i$ this normalization is the displayed smooth family $U_i\times\mathbb C_t$. All its omitted sets are nowhere dense analytic subsets.

The sheaf $\mathcal F=\nu_*\mathcal O_Y$ is a finite coherent $\mathcal O_C$-module. We spell out the finite-map point here. The root cover is a closed subspace of the monic cover $b^q=a_i g_i^{q-1}$, whose algebra is finite free over $\mathcal O_C$. Near its finite fiber, the normalization construction in \cref{thm:analytic-normalization} uses finitely many coordinates integral over the root-cover local rings. These coordinates are integral over $\mathcal O_C$ as well. Indeed the algebra generated by $b$ and any finite list of those coordinates is a finite base module: reduce powers of $b$, and then powers of each new coordinate by its monic equation. The determinant argument for multiplication on this finite module gives a monic base equation for each coordinate. Apply this at all points of the finite fiber and shrink to representatives of the equations.

Here is why this gives coherence on a neighborhood, rather than only finite generation of a stalk. Embed the finitely many source neighborhoods using these integral coordinates. They are closed analytic subspaces of successive monic covers of a base neighborhood. Lift the monic equations to an ambient neighborhood of $C$; its coherent defining ideal restores $C$ after division. For one coordinate, separate the root clusters by contour integrals and use their monic factors. The argument in \cref{lem:strong-analytic-noetherian} reduces a finite presentation modulo such a factor and represents coordinate multiplication by its companion matrix. It gives a finite matrix of holomorphic base functions presenting the direct image, and makes that direct image exact. Repeat for each coordinate. Coherence of the analytic defining ideals and exactness give a finite presentation for the closed source subspaces too. Properness confines all nearby fibers to the chosen source neighborhoods, so their finite direct sum describes $\mathcal F$. This proves the assertion using only monic division and Oka coherence. Bounded removal on the singular base is not needed.

There is a holomorphic action of $\mathbb C^*$ on $C'$ given by
\[
 \lambda\cdot(z,a,b)=(z,\lambda^q a,\lambda b).
\]
It lifts uniquely to $Y$. Indeed a product of a normal space with a disc is normal: a bounded holomorphic function off an analytic subset first extends on the regular product locus, and its Cauchy coefficients in the disc variable extend on the normal factor by \cref{lem:normal-locus-open}; the locally uniformly convergent Cauchy series gives the extension on the product. Apply the normalization factorization in \cref{cor:finite-bimeromorphic-normal-target} to $\mathbb C^*\times Y\to C'$. Its inverse image of the singular locus is nowhere dense. Uniqueness gives the action law. On the dense open family this action is $t\mapsto\lambda t$.

Use an $S^1$-invariant neighborhood of $(x,0)$ in $C$. The finite-dimensional vector space $\mathcal F_{(x,0)}/\mathfrak m_{(x,0)}\mathcal F_{(x,0)}$ carries the circle action. Averaging any Hermitian inner product makes it unitary; commuting unitary matrices have a common eigenbasis. Each continuous one-dimensional circle character is $\lambda\mapsto\lambda^d$ for an integer $d$, as follows by lifting its argument on a small interval and using the additive character equation. Lift this eigenbasis to holomorphic sections and project each lift to its character by
\[
 h_d(y)=\frac1{2\pi}\int_0^{2\pi}
         \mathrm e^{-d\mathrm i\theta}
         h(\mathrm e^{\mathrm i\theta}\cdot y)\,\mathrm d\theta.
\]
On the regular locus differentiation under this compact integral proves holomorphy. Its local bound and normal removal extend it across the singular locus. It projects to the chosen eigenvector. These finitely many homogeneous sections $h_j$, of weights $d_j$, generate $\mathcal F$ near $(x,0)$ by the local-ring determinant argument in \cref{lem:normal-locus-open}. After shrinking the base neighborhood, they generate on an invariant neighborhood of its zero section. Every $d_j$ is nonnegative: on the dense open family a homogeneous function is $f_j(z)t^{d_j}$, and holomorphy at $t=0$ excludes a negative weight. The same identity follows for nonzero $t$ from the circle character equation and the one-variable Laurent expansion.

Consider a section $h$ of weight $p$. Write it locally as $\sum_j k_jh_j$, with $k_j$ holomorphic on $C$, and project this equation to weight $p$. The resulting coefficient of $h_j$ has weight $p-d_j$. A holomorphic coefficient on $C$ is the restriction of an ambient holomorphic function near $(x,0)$. Projecting its convergent power series under $a\mapsto\lambda^q a$ shows that it is a homogeneous polynomial in $a$, with holomorphic coefficients in $z$, of degree $(p-d_j)/q$; it is zero unless this degree is a nonnegative integer.

\Needspace{4\baselineskip}
Consequently the weight-$p$ sections are generated over $\mathcal O_U$ by the finite list
\[
 a^\alpha h_j,\qquad q|\alpha|+d_j=p.
\]
The statement is local at $(x,0)$; shrinking once gives representatives of all these generators on a common invariant neighborhood.

We identify that list with generators of the desired ideal. On $\{g_i\ne0\}$ every weight-$p$ section has the form $f(z)t^p$. Its coefficient $f$ is holomorphic there. Choose a smaller base neighborhood and a fixed sufficiently small $\delta>0$. For such a $z$, evaluate at any $t$ with
\[
 |t|^q|g(z)|=\delta.
\]
The images $(z,a)$ stay in a fixed compact subset of the chosen neighborhood of $C$, since $|a|=\delta$. Finite properness of $\nu$ makes their inverse images compact. Thus $h$ is bounded there and
\[
 |f(z)|\leq C\delta^{-p/q}|g(z)|^{p/q}.
\]
Bounded holomorphic removal across $\{g_i=0\}$ extends $f$ to $U$, and continuity preserves this inequality. Each of the finite homogeneous generators therefore gives a holomorphic member of $\mathcal J_{p/q}(g)$.

Conversely, if $f\in\mathcal J_{p/q}(g)$, the function $f(z)t^p$ on the dense open family satisfies
\[
 |f(z)t^p|\leq C|t^qg(z)|^{p/q}=C|a|^{p/q}.
\]
It is locally bounded on $Y$ and extends holomorphically there by normality, \cref{lem:normal-locus-open}. It remains homogeneous by density. The preceding finite homogeneous generation then writes its coefficient $f$ as a holomorphic base combination of the finite list of coefficient functions. Thus this list generates $\mathcal J_{p/q}(g)$ near $x$. Oka coherence proves that this ideal is coherent.

Finally a local analytic presentation $v=(p/q)\log\sum|g_j|^2+O(1)$ identifies $\mathcal I_\infty(v)$ with this ideal. The definitions agree on overlaps, so local coherence gives the global assertion. A nonvanishing tuple gives the unit ideal. In dimension zero every finite weight likewise gives the unit ideal.
\end{proof}
\index{normalization!bounded growth ideals}\index{coherence!bounded growth ideals}

\begin{corollary}\label{cor:analytic-integrability-bounded-growth}
If $v$ has analytic singularities on compact $X$, then for every $\varepsilon>0$ there is $k_0$ such that
\[
 \mathcal I(k(1+\varepsilon)v)\subseteq\mathcal I_\infty(kv)
 \quad(k\geq k_0).
\]
\end{corollary}
\begin{proof}
Apply \cref{lem:analytic-weighted-integrability-growth} with $t=k(1+\varepsilon)$ and $\eta=\varepsilon/(1+\varepsilon)$. Its finite cover gives one threshold for all the germs, and $t(1-\eta)=k$. The bounded-growth definition gives the inclusion. Dimension zero is immediate.
\end{proof}
